\documentclass[10pt,a4paper]{amsart}
\usepackage[margin=19mm]{geometry}
\usepackage{amsmath,amssymb,mathtools}
\usepackage[scr=rsfs,cal=euler]{mathalfa}
\usepackage{xfrac}
\usepackage[unicode,hidelinks,bookmarks=false]{hyperref}
\newcommand{\ie}{i.e.\ }
\newcommand{\cf}{cf.\ }
\newcommand{\resp}{respectively\ }
\newcommand{\Subsection}{\subsection}
\providecommand{\nobreakdash}{\nobreak\hbox{-}}
\setlength{\emergencystretch}{2em}
\multlinegap0pt
\usepackage{bm}
\usepackage{bbm}
\usepackage{dsfont}
\usepackage{xspace}
\usepackage{cancel}
\usepackage{mathtools}
\usepackage{amsthm}
\makeatletter
\@ifundefined{theorem}{\newtheorem{theorem}{Theorem}}{}
\@ifundefined{lemma}{\newtheorem{lemma}[theorem]{Lemma}}{}
\@ifundefined{prop}{\newtheorem{prop}[theorem]{Proposition}}{}
\makeatother
\title[Random partition for Tokushige's $r$-wise intersecting conje]{Random partition for Tokushige's $r$-wise intersecting conjecture}
\author{Yongjiang Wu, Lihua Feng}
\date{Source: arXiv:2606.31075v1 (CC BY 4.0)}
\begin{document}
\maketitle

\noindent\emph{Source and licence.} Random partition for Tokushige's $r$-wise intersecting conjecture. Version arXiv:2606.31075v1 (CC BY 4.0), distributed under the Creative Commons Attribution 4.0 International licence, as recorded in the arXiv licence metadata for this exact version and shown on the abstract page https://arxiv.org/abs/2606.31075v1. Full rights notice: licenses/arxiv-2606.31075-CC-BY-4.0.txt.\par\smallskip

\noindent\textit{Editorial scope.} Counted unit: the Prop whose source label is \texttt{prop:slicebound} in the article (source0.tex lines 338--345, source label \texttt{prop:slicebound}), delivered together with the article's own complete original proof of it (source0.tex lines 347--450, one \texttt{proof} environment of 104 lines). No mathematics is written, repaired or re-derived: statement and proof are the article's own text line for line. Required context retained uncounted: eq:slice-measure (lines 257--259); lem:slice (lines 261--267); lem:partition (lines 281--288). Every definition, display and label of those lines is retained unchanged. Omitted from this package and not redistributed: 1--256, 260--260, 268--280, 289--337, 346--346, 451--701. Nothing inside a retained excerpt is omitted or altered: every retained body is delivered exactly as the article's lines are written, so no omission receipt of any kind accompanies this package. Citations: the retained excerpts carry no citation command at all, so no bibliography entry of the article is transcribed and no reference is redistributed. Reference adaptation: none was needed and none was made; every reference inside the delivered unit resolves inside it, so no unresolved reference or citation is produced. Printed slips kept literal: no printed word, symbol, hypothesis or number of the counted statement or of its proof was altered. Layout and class: the article's own class is \texttt{article} with packages including amsmath, amsfonts, amsthm, amssymb, mathtools, mathrsfs, graphicx, color, latexsym, tikz, calc, tikz-cd, epstopdf, enumerate; the wrapper is the lane's pinned \texttt{amsart} editorial wrapper at 10pt on A4, which restates the theorem-like environments the retained text needs and adds the article's own macro definitions that the retained text needs (none), with extra wrapper packages (bm, bbm, dsfont, xspace, cancel, mathtools). The delivered numbering is the unit's own. Assets: no retained span contains a figure environment, a graphics-inclusion command, a PDF-page inclusion, a table or a TikZ picture; no image file is copied into this package, the package declares no asset dependency and no image file is redistributed with it. Licence: version arXiv:2606.31075v1 of this article is distributed under the Creative Commons Attribution 4.0 International licence, as recorded in the arXiv licence metadata for this exact version and shown on the abstract page https://arxiv.org/abs/2606.31075v1. A full rights notice is delivered at licenses/arxiv-2606.31075-CC-BY-4.0.txt. Characters: every retained excerpt is pure ASCII, so the delivered engine, XeLaTeX with its default Latin Modern text font, reports no missing character.
\begin{equation}\label{eq:slice-measure}
 \mu(\mathcal{A})=\sum_{B\subseteq H}p^{|B|}q^{m-|B|}w_B.
\end{equation}

\begin{lemma}\label{lem:slice}
If $B_1,\ldots,B_r\subseteq H$ and
$B_1\cap\cdots\cap B_r=\emptyset$, then
$$
 w_{B_1}+\cdots+w_{B_r}\le r-1.
$$
\end{lemma}

\begin{lemma}[Partition law]\label{lem:partition}
Let $S$ be a finite set with $1\leq |S|=k\le s$ and let $0<t\le\frac{1}{s}$.  There is a
probability distribution on set partitions $\Pi$ of $S$ (with nonempty blocks only) such that, for every
nonempty $D\subseteq S$,
$$
 \mathbb P(D\in\Pi)=t^{|D|-1}(1-t)^{k-|D|}.
$$
\end{lemma}

\begin{prop}[Supercritical slice bound]\label{prop:slicebound}
In the setting of this section, we have
$$
 \mu(\mathcal{A})\le p-p\left(1-(r-1)\frac{q}{p}\right)(1-w_H).
$$
In particular $\mu(\mathcal{A})\le p$.  Moreover, equality $\mu(\mathcal{A})=p$ forces
$w_H=1$, equivalently $\mathcal{A}_H=2^T$.
\end{prop}

\begin{proof}
Fix a coordinate $1\in H$ and put $U=H\setminus\{1\}$. Since  $|H|=m\le r$, we have $|U|=m-1\le r-1=s$.
For $B,D\subseteq U$, define
$$
 x_B=w_B,
 \qquad y_D=1-w_{H\setminus D}.
$$
We first prove a local inequality.  Fix $B\subseteq U$.  If
$D_1,\ldots,D_s\subseteq B$ and $D_1\cup\cdots\cup D_s=B$, then
$$
 B\cap(H\setminus D_1)\cap\cdots\cap(H\setminus D_s)=\emptyset.
$$
By Lemma \ref{lem:slice},
\begin{equation}\label{eq:local-cover}
 x_B=w_B\le\sum_{j=1}^s\bigl(1-w_{H\setminus D_j}\bigr)
        =\sum_{j=1}^s y_{D_j}.
\end{equation}

For $B\ne\emptyset$, choose the random partition $\Pi$ of $B$ from
Lemma \ref{lem:partition} and add $s-|\Pi|$ empty blocks.  Taking expectations in
\eqref{eq:local-cover}, and using
$$
\sum_{j=1}^s y_{D_j}
=
\sum_{\emptyset\neq D\subseteq B,\ D\in\Pi} y_D
+
(s-|\Pi|)\,y_\emptyset,\quad \mathbb E[|\Pi|]=\sum_{\emptyset\ne D\subseteq B}\mathbb P(D\in\Pi)
          =\frac{1-(1-t)^{|B|}}{t},
$$
we obtain
\begin{align}
 x_B&\le
\mathbb E\left[\sum_{\emptyset\neq D\subseteq B}\mathbf 1_{\{D\in\Pi\}}\,y_D\right]
+
\mathbb E[s-|\Pi|]\,y_\emptyset\notag\\
 &=\sum_{\emptyset\ne D\subseteq B}t^{|D|-1}(1-t)^{|B|-|D|}y_D
     +\left(s-\frac{1-(1-t)^{|B|}}{t}\right)y_\emptyset \notag\\
 &=\sum_{D\subseteq B}t^{-1}t^{|D|}(1-t)^{|B|-|D|}y_D
     -\frac{1-st}{t}y_\emptyset .        \label{eq:local-ineq}
\end{align}

For $B=\emptyset$, take $D_1=\cdots=D_s=\emptyset$. These sets cover $B$. Applying the local inequality \eqref{eq:local-cover} to this cover gives
$
x_\emptyset \le s y_\emptyset,
$
which is precisely \eqref{eq:local-ineq} in the case $B=\emptyset$





Multiply \eqref{eq:local-ineq} by $t p^{|B|}q^{|U|-|B|}$ and sum over
$B\subseteq U$. The left-hand side becomes
$$
t\sum_{B\subseteq U}p^{|B|}q^{|U|-|B|}x_B.
$$
The right-hand side becomes
\begin{align*}
&\sum_{B\subseteq U}\sum_{D\subseteq B}
p^{|B|}q^{|U|-|B|}t^{|D|}(1-t)^{|B|-|D|}y_D-(1-st)y_\emptyset\\
=&\sum_{D\subseteq U}y_D
\sum_{B\supseteq D}
p^{|B|}q^{|U|-|B|}t^{|D|}(1-t)^{|B|-|D|}-(1-st)y_\emptyset.
\end{align*}
 For fixed $D\subseteq U$, write $B=D\cup E$ with $E\subseteq U\setminus D$. Then
\begin{align*}
 &\sum_{B\supseteq D}
p^{|B|}q^{|U|-|B|}t^{|D|}(1-t)^{|B|-|D|}
=
p^{|D|}t^{|D|}
\sum_{E\subseteq U\setminus D}
p^{|E|}q^{|U|-|D|-|E|}(1-t)^{|E|}\\
=&
p^{|D|}t^{|D|}
\bigl(p(1-t)+q\bigr)^{|U|-|D|}
=
p^{|D|}t^{|D|}
p^{|U|-|D|}
=
q^{|D|}p^{|U|-|D|},
\end{align*}
where we use the binomial theorem, $p(1-t)+q=p$, and $pt=q$.
Therefore, 
\begin{equation}\label{eq:weighted-local}
 t\sum_{B\subseteq U}p^{|B|}q^{|U|-|B|}x_B
 \le \sum_{D\subseteq U}q^{|D|}p^{|U|-|D|}y_D-(1-st)y_\emptyset.
\end{equation}

Now compare $\mathcal{A}$ with the star centred at coordinate $1$.  Using
\eqref{eq:slice-measure}, we obtain
\begin{align*}
 \mu(\mathcal{A})-p
 &=q\sum_{B\subseteq U}p^{|B|}q^{|U|-|B|}x_B
   -p\sum_{D\subseteq U}q^{|D|}p^{|U|-|D|}y_D \\
 &=p\left(
 t\sum_{B\subseteq U}p^{|B|}q^{|U|-|B|}x_B
 -\sum_{D\subseteq U}q^{|D|}p^{|U|-|D|}y_D\right) \\
 &\le -p(1-st)y_\emptyset.
\end{align*}
Since $y_\emptyset=1-w_H$ and $pst=(r-1)q$, the inequality follows.
Observe that $1-st>0$.
If $\mu(\mathcal{A})=p$, then $w_H=1$. Thus, $\mathcal{A}_H$ has tail measure one. Since every tail
atom has positive measure, this is equivalent to $\mathcal{A}_H=2^T$.
\end{proof}

\end{document}
